\section{Abstract}

This thesis studies price tracking in a synthetic, zero-fee constant-product automated market maker (AMM) against an exogenous centralized-exchange (CEX) reference. It separates the swap rule's algebraic implications from effects that depend on agent behavior. In an isolated pool, an arbitrageur limited to a fraction $c$ of the input reserve can reduce out-of-band log error by at most $2\log(1+c)$ per tick. For a log shock of $0.20$ and a band of $0.003$, the idealized time to the band is 20 ticks at $c=0.005$, four at $c=0.03$, and two at $c=0.10$. This calculation explains the direction of the configured arbitrage result; simulation measures departures under noise and other rules.

Across 300 Monte Carlo seeds and 13,500 runs, the high-capacity scenario has a mean excess-error area of $0.10$, compared with $2.36$ without arbitrage. The result is conditional on the specified rule and parameter grid, and the sign of the effect is mechanically built into that rule. The current liquidity-withdrawal trigger fires in every adaptive-condition run, usually at the first tick. Its apparent increase in peak response occurs only at the shallowest depth with fixed absolute orders; it disappears when orders scale with reserves. The original peak-displacement metric is not overshoot relative to the new reference, so this sensitivity does not establish a general liquidity-amplification effect.

The more transferable simulator finding is about persistence: 31.2\,\% of paths that returned to the no-arbitrage band later left it; at depth factor $0.5$, the share was 98.8\,\%. A first return should therefore be distinguished from sustained tracking when defining monitoring measures. The two-state reduced map improves on persistence in some conditions but does not outperform a scalar AR(1). A band-aware map improves recursive replay error by about 1.7\,\%, concentrated in one condition and using the realized future CEX path; it is not a real-time forecasting result. The contribution is a reproducible, mechanism-focused simulation study that distinguishes algebraic results from rule-dependent sensitivities and reports where its reduced models fail. Event replay with synchronized market data and realistic costs remains future work.

\textbf{Keywords:} automated market maker; arbitrage; price tracking; liquidity withdrawal; agent-based simulation; Monte Carlo experiment
